\documentclass[11pt]{article}
\usepackage[a4paper, margin=2.5cm]{geometry}
\usepackage{amsmath, amssymb, amsthm}
\usepackage{enumitem}
\usepackage[hidelinks]{hyperref}

\newtheorem{theorem}{Theorem}
\newtheorem{lemma}[theorem]{Lemma}
\newtheorem{proposition}[theorem]{Proposition}
\theoremstyle{definition}
\newtheorem{definition}{Definition}
\newtheorem{example}{Example}
\theoremstyle{remark}
\newtheorem*{remark}{Remark}

\title{Course Notes: An Example in Analysis}
\author{Your Name}
\date{\today}

\begin{document}
\maketitle

\begin{abstract}
These notes demonstrate a theorem-friendly template: numbered theorem
environments, clean proofs, and display math — ideal for course notes and
problem sets.
\end{abstract}

\section{Definitions and Theorems}

\begin{definition}[Continuity]
A function $f : \mathbb{R} \to \mathbb{R}$ is continuous at $x_0$ if
\[
\forall \varepsilon > 0 \;\; \exists \delta > 0 :
|x - x_0| < \delta \implies |f(x) - f(x_0)| < \varepsilon .
\]
\end{definition}

\begin{theorem}[Intermediate Value]
Let $f$ be continuous on $[a,b]$. If $f(a) \cdot f(b) < 0$, then there
exists $c \in (a,b)$ with $f(c) = 0$.
\end{theorem}

\begin{proof}
WLOG assume $f(a) < 0 < f(b)$. Let
\[
S = \{ x \in [a,b] : f(x) < 0 \},
\]
and take $c = \sup S$. By continuity, $f(c) = 0$.
\end{proof}

\begin{remark}
The converse fails: a function may have a root without changing sign.
\end{remark}

\section{Exercises}
\begin{enumerate}[label=\textbf{E\arabic*.}]
  \item Prove that a monotone function has at most countably many
        discontinuities.
  \item Show that $f(x) = x\sin(1/x)$ (extended by $f(0)=0$) is continuous
        on $\mathbb{R}$.
\end{enumerate}

\end{document}
